\section{The \pwm model}\label{sec:model}

\paragraph{Station tokens.}
Each station becomes one token, built by a two-layer MLP from:
\begin{itemize}
\item its last 72 hours of temperature, humidity and wind, with observed flags;
\item its last 24 hours of the other channels (wind direction, dew point or solar radiation, pressure, visibility or
PM$_{2.5}$, depending on the league);
\item its mean daily cycle over the last 14 days and its trailing 7-day level;
\item the sine of the sun's elevation at the station for each of the 24 forecast hours;
\item its position, at a fixed scale that is the same in every network, and the day of year.
\end{itemize}
Hour-of-day and month embeddings are added. There is no per-station embedding, so the same weights run on any
station network.

\paragraph{Mixing across stations.}
Three blocks of multi-scale graph attention mix the tokens. Each head's attention bias is a learned mixture of the
$k$-nearest-neighbour graph, its square and the identity, so different heads specialise in different ranges. Region
tokens, pooled over spectral clusters, and a global token carry information between distant stations. The model has
width 128 and 4 heads.

\paragraph{Output.}
A linear head gives the 10\%, 50\% and 90\% quantiles for each variable and each of the 24 hours. They are
expressed as offsets, in units of the network's spread, from a base: the recent daily cycle plus the latest anomaly.
Training minimises the quantile (pinball) loss \citep{koenker1978quantiles}. Training randomly darkens stations and
blanks bursts of hours before the origin, so that the model learns to handle real outages. The reported forecast is
the mean of three seeds. Intervals are calibrated on validation by conformalised quantile regression
\citep{romano2019cqr}, but this paper scores only the median.

\paragraph{Hybrid.}
The hybrid also reads global forecasts, using for each model the newest run issued before the forecast starts: the
1-day-old run for hours 0--16 and the 2-day-old run after. Each forecast is bias-corrected by its trailing 30-day
error at that station and hour of day. Two hybrid designs compete on validation:
\begin{itemize}
\item \textbf{H1} adds the corrected forecasts of IFS 0.25, AIFS, ICON and GFS as inputs;
\item \textbf{H2} replaces the base with their mean, a model-output-statistics base, and learns only the correction.
\end{itemize}

\paragraph{Additions in later versions (\S\ref{sec:versions}).}
\begin{itemize}
\item \textbf{Cross-network training (v0.6).} Half of the training batches come from another station network. It uses only
that network's hours before this league's validation period, with its own normalisation statistics from those hours
alone. Other channels are matched by name, and the clock is shifted to the league's.
\item \textbf{14-day encoder (v0.8).} The full 14 $\times$ 24-hour history, previously only averaged into a daily
cycle, passes day by day through a small MLP with attention pooling over days.
\item \textbf{Weight averaging (v0.8).} An exponential moving average of the weights (decay 0.998) is what gets
validated and used.
\end{itemize}
The final models have 0.93--1.13M parameters.
